\chapter{Tổng quan và bối cảnh nghiên cứu}
\label{chap:tong_quan}

\section{Tổng quan về hệ thống giao thông thông minh (ITS)}
Hệ thống giao thông thông minh (Intelligent Transportation Systems -- ITS) là sự kết hợp giữa công nghệ thông tin, truyền thông và các thuật toán điều khiển hiện đại nhằm nâng cao hiệu quả vận hành mạng lưới giao thông, bảo đảm an toàn, tối ưu hóa năng lượng và giảm thiểu ô nhiễm môi trường đô thị.

Theo các khung kiến trúc chuẩn mực quốc tế (như ISO 14813), một hệ thống ITS đô thị thường gồm ba phân hệ chính:
\begin{enumerate}
    \item \textbf{Hệ thống quản lý giao thông nâng cao (ATMS):} Giám sát trạng thái mạng lưới đường bộ thời gian thực qua hệ thống cảm biến, phát hiện sự cố tự động, điều khiển đèn tín hiệu thích ứng và điều phối lực lượng xử lý tình huống phát sinh.
    \item \textbf{Hệ thống thông tin người tham gia giao thông (ATIS):} Cung cấp thông tin cập nhật liên tục về tình trạng lưu thông, thời gian hành trình dự kiến và hướng dẫn lộ trình tối ưu thông qua biển báo điện tử VMS, ứng dụng di động và đài phát thanh giao thông.
    \item \textbf{Hệ thống vận tải công cộng nâng cao (APTS):} Hỗ trợ quản lý và tối ưu biểu đồ chạy xe buýt, metro, cung cấp thời gian đến bến dự kiến (ETA) và điều phối ưu tiên qua nút giao cho phương tiện công cộng.
\end{enumerate}

Trong đó, bài toán \textbf{dự báo dòng giao thông đô thị (urban traffic forecasting)} đóng vai trò trung tâm trong phân hệ ATMS và ATIS. Việc dự báo chính xác trạng thái mạng lưới trước 15 đến 30 phút cho phép các nhà quản lý chuyển từ trạng thái phản ứng bị động sang điều hành chủ động, kịp thời điều chỉnh chu kỳ đèn tín hiệu và phân luồng từ xa trước khi ùn tắc lan rộng.

\subsection{Tiến trình phát triển của hệ thống giao thông thông minh}
Quá trình phát triển của ITS trải qua ba giai đoạn công nghệ chính:
\begin{enumerate}
    \item \textbf{Thế hệ ITS 1.0 (1970 -- 1990): Kỷ nguyên cảm biến cơ điện cố định.}
    Giai đoạn này gắn liền với các cảm biến vòng từ (inductive loop detectors) chôn dưới mặt đường và hệ thống đèn tín hiệu phối hợp cố định theo thời gian trong ngày (như phương pháp Webster). Dữ liệu thu được chủ yếu phục vụ các tuyến cao tốc, thu thập lưu lượng tổng hợp trên từng làn và chưa có khả năng liên kết mạng lưới không gian cũng như phân loại phương tiện.
    \item \textbf{Thế hệ ITS 2.0 (1990 -- 2015): Kỷ nguyên truyền thông không dây và định vị vệ tinh.}
    Sự phổ biến của hệ thống GPS và mạng di động mở ra nguồn dữ liệu từ phương tiện di động (Floating Car Data -- FCD). Các hệ thống điều khiển đèn thích ứng như SCOOT (Anh) và SCATS (Úc) bắt đầu ứng dụng mô hình hồi quy để điều hòa lưu lượng. Tuy nhiên, chi phí đầu tư và bản quyền của các giải pháp này rất lớn, chủ yếu thiết kế cho dòng xe ô tô chiếm ưu thế tại các quốc gia phát triển.
    \item \textbf{Thế hệ ITS 3.0 (Từ 2015 đến nay): Kỷ nguyên AIoT, thị giác máy tính và học sâu đồ thị.}
    Sự phát triển mạnh mẽ của thị giác máy tính và mạng nơ-ron học sâu giúp tận dụng trực tiếp hạ tầng camera giám sát công cộng (CCTV) sẵn có tại các đô thị. Dữ liệu hình ảnh được chuyển hóa thành các luồng số liệu giao thông liên tục, kết hợp với các mô hình mạng nơ-ron đồ thị không -- thời gian (STGNN) để dự báo toàn mạng lưới theo thời gian thực với chi phí phần cứng hợp lý.
\end{enumerate}

\section{Động học dòng giao thông đô thị và bài toán dự báo đa bước}
\subsection{Đặc tính động học phi tuyến và tính phụ thuộc không -- thời gian}
Dòng giao thông đô thị là một hệ thống động học phi tuyến phức tạp, chịu sự chi phối đồng thời của hai yếu tố tương quan:
\begin{enumerate}
    \item \textbf{Tương quan thời gian (temporal dependencies):} Lưu lượng xe mang tính quy luật rõ rệt theo chu kỳ ngày đêm và chu kỳ tuần (giữa ngày làm việc và ngày nghỉ), đồng thời có quán tính ngắn hạn. Khi dòng phương tiện đông đúc xuất hiện tại một thời điểm, trạng thái này có xu hướng duy trì trong các khoảng thời gian kế tiếp trước khi giải tỏa.
    \item \textbf{Tương quan không gian (spatial dependencies):} Giao thông vận hành trên mạng lưới đường phố liên thông có cấu trúc đồ thị. Trạng thái lưu thông tại một nút giao phụ thuộc trực tiếp vào dòng xe từ các nút thượng lưu đổ về và khả năng thoát xe tại các nút hạ lưu. Mối tương quan này bị ràng buộc bởi hình học mạng lưới, khoảng cách giữa các giao lộ, quy định đường một chiều và số làn xe.
\end{enumerate}

\subsection{Đặc thù dòng giao thông hỗn hợp và động học di chuyển của xe máy}
Tại các đô thị ở Việt Nam, dòng giao thông mang những đặc thù khác biệt đáng kể so với mô hình giao thông tại các nước phát triển:
\begin{itemize}
    \item \textbf{Dòng chảy phi làn (non-lane based flow):} Khác với giả định phương tiện di chuyển bám làn thành hàng dọc trong các lý thuyết giao thông cổ điển, xe máy tại Việt Nam di chuyển linh hoạt theo dòng chảy tự do, liên tục tận dụng các khoảng trống giữa các xe ô tô để vượt lên.
    \item \textbf{Đặc thù di chuyển theo cụm mật độ cao:} Khi dừng chờ đèn đỏ tại các ngã tư, hàng trăm xe máy dồn ứ lại thành cụm dày đặc phía trước ô tô. Khi đèn chuyển xanh, cụm xe này đồng loạt bứt tốc với gia tốc lớn, tạo ra dao động lưu lượng đột biến trong chu kỳ quan sát 5 phút.
    \item \textbf{Hệ số xe con quy đổi biến thiên:} Trong điều kiện lưu thông bình thường, một xe máy tương đương khoảng 0.2 -- 0.25 xe con quy đổi (PCU). Tuy nhiên trong điều kiện ùn ứ tốc độ thấp, xe máy luồn lách dễ dàng hơn trong khi ô tô chiếm dụng diện tích mặt đường lớn. Do đó, việc tách biệt chính xác số lượng ô tô và xe máy là yêu cầu quan trọng để phản ánh đúng hiện trạng tải trọng đường bộ.
\end{itemize}

\subsection{Ý nghĩa của tầm nhìn dự báo 15--30 phút}
Bài toán dự báo dòng giao thông đa bước đặt mục tiêu ước lượng lưu lượng phương tiện qua $H$ bước thời gian liên tiếp trong tương lai:
\begin{equation}
    \hat{\mathbf{Y}}_{t+1}, \hat{\mathbf{Y}}_{t+2}, \dots, \hat{\mathbf{Y}}_{t+H}
\end{equation}
với chu kỳ lấy mẫu tiêu chuẩn $\Delta t = 5$ phút, tầm nhìn dự báo $H = 6$ bước tương ứng với 30 phút tới.

Tầm nhìn dự báo 15--30 phút có ý nghĩa thực tiễn quan trọng:
\begin{itemize}
    \item \textbf{Thời gian đáp ứng của chu kỳ đèn tín hiệu:} Các giải pháp điều khiển đèn phối hợp cần khoảng thời gian tương đương 3 đến 5 chu kỳ đèn để chuyển đổi kế hoạch pha mượt mà mà không gây gián đoạn dòng xe. Dự báo trước 15--30 phút giúp trung tâm điều hành kịp thời điều chỉnh thời lượng đèn xanh đón đầu dòng xe đang tiếp cận.
    \item \textbf{Mô hình lan truyền sóng tắc nghẽn:} Theo mô hình thủy động lực học giao thông (LWR), khi lưu lượng vượt quá dung lượng mặt đường, sóng tắc nghẽn sẽ lan truyền ngược chiều dòng xe với vận tốc khoảng 15--25 km/h. Dự báo sớm 30 phút cho phép hệ thống kịp thời thông báo và điều tiết từ xa, hỗ trợ phân luồng phương tiện trước khi điểm nghẽn hình thành.
\end{itemize}

\section{Khảo cứu các kỹ thuật nhận dạng và đếm phương tiện từ camera}
Để cung cấp số liệu lưu lượng cho các mô hình dự báo mạng lưới, khâu thị giác máy tính đóng vai trò thu thập tín hiệu đầu vào. Hiện nay có ba hướng tiếp cận chính:

\subsection{Tiếp cận bám vết theo phát hiện (tracking-by-detection)}
Phương pháp này kết hợp mô hình phát hiện vật thể (như YOLO, Faster R-CNN) với thuật toán bám vết liên khung hình (như ByteTrack, DeepSORT).
\begin{itemize}
    \item \textbf{Nguyên lý:} Tại mỗi khung hình, mô hình phát hiện xác định các bounding box của từng xe. Thuật toán bám vết liên kết quỹ đạo phương tiện qua các khung hình liên tiếp dựa trên bộ lọc Kalman và chỉ số IoU (Intersection-over-Union). Khi một quỹ đạo cắt qua vạch ảo định sẵn, hệ thống ghi nhận một lượt xe.
    \item \textbf{Hạn chế trong giao thông xe máy mật độ cao:}
    \begin{enumerate}
        \item \textit{Hiện tượng che khuất lẫn nhau:} Vào giờ cao điểm, các xe máy đi sát nhau khiến tỷ lệ che khuất thân xe rất lớn, dẫn đến việc mô hình phát hiện bỏ sót hoặc sinh ra các hộp nhận dạng sai lệch.
        \item \textit{Chi phí tính toán của khâu NMS:} Khi số lượng xe trong khung hình lên tới hàng trăm, mô hình sinh ra hàng nghìn bounding box ứng viên. Thuật toán triệt tiêu cực đại cục bộ (NMS) với độ phức tạp $\mathcal{O}(M^2)$ trở thành điểm nghẽn tính toán trên các thiết bị biên.
        \item \textit{Gián đoạn thời gian trên camera công cộng:} Để tiết kiệm băng thông hạ tầng mạng, nhiều camera giao thông chỉ truyền ảnh tĩnh định kỳ (ví dụ 60 giây/ảnh), làm mất tính liên tục khung hình và khiến các thuật toán bám vết mất dấu chuyển động.
    \end{enumerate}
\end{itemize}

\subsection{Tiếp cận ước lượng bản đồ mật độ (density map estimation)}
Xuất phát từ bài toán đếm đám đông (crowd counting), các mô hình như CSRNet, CAN chuyển đổi bài toán đếm rời rạc thành bài toán hồi quy bản đồ mật độ liên tục:
\begin{itemize}
    \item \textbf{Nguyên lý:} Mỗi phương tiện được biểu diễn bằng một điểm tâm, bản đồ mật độ được tạo ra qua phép tích chập với hàm Gaussian. Tổng số lượng xe được tính bằng tích phân trên toàn bộ bức ảnh.
    \item \textbf{Hạn chế:} Bản đồ mật độ không phân biệt được cụ thể loại phương tiện (ô tô hay xe máy). Trong khi đó, việc mất thông tin phân lớp gây khó khăn cho việc đánh giá tải trọng giao thông. Ngoài ra, việc gán nhãn điểm tâm cho tập dữ liệu lớn đòi hỏi chi phí nhân công đáng kể.
\end{itemize}

\subsection{Tiếp cận hồi quy đếm xe trực tiếp (direct regression)}
Nhằm khắc phục các nhược điểm trên, đề tài lựa chọn và phát triển hướng tiếp cận \textbf{hồi quy đếm xe trực tiếp theo phân lớp}.
\begin{itemize}
    \item \textbf{Nguyên lý:} Mô hình học sâu nhận đầu vào là ảnh tĩnh $\mathbf{I} \in \mathbb{R}^{H \times W \times 3}$ và ánh xạ trực tiếp sang vector số lượng hai chiều:
    \begin{equation}
        f_\theta(I) \longrightarrow \hat{\mathbf{c}} = \begin{bmatrix} \hat{c}_{\text{car}} \\ \hat{c}_{\text{bike}} \end{bmatrix} \in \mathbb{R}^2
    \end{equation}
    \item \textbf{Ưu điểm:}
    \begin{enumerate}
        \item \textit{Không cần khâu NMS:} Mô hình không sinh bounding box, loại bỏ độ phức tạp tính toán bậc hai và đạt tốc độ xử lý nhanh (dưới 35 ms/ảnh trên phần cứng phổ thông).
        \item \textit{Bảo toàn phân lớp phương tiện:} Tách biệt rõ ràng số lượng ô tô và xe máy, cung cấp đầu vào chính xác cho mô hình dự báo không -- thời gian.
        \item \textit{Phù hợp với luồng ảnh tĩnh định kỳ:} Xử lý độc lập trên từng ảnh, không đòi hỏi video liên tục băng thông cao.
        \item \textit{Khắc phục hiện tượng học bám cảnh nền:} Áp dụng các kỹ thuật tăng cường dữ liệu ngẫu nhiên (Color Jittering, Random Erasing) và kiểm chứng sự tập trung đặc trưng qua bản đồ nhiệt Grad-CAM++.
    \end{enumerate}
\end{itemize}

\section{Khảo cứu các mô hình dự báo chuỗi thời gian giao thông}
Bài toán dự báo dòng giao thông không -- thời gian đã trải qua các giai đoạn phát triển chính:

\subsection{Các mô hình thống kê truyền thống}
Trong giai đoạn đầu, các phương pháp thống kê cổ điển giữ vai trò chủ đạo:
\begin{itemize}
    \item \textbf{ARIMA và SARIMA:} Phân tích chuỗi thời gian dựa trên các thành phần tự hồi quy, sai phân và trung bình trượt dưới giả định chuỗi dừng và tuyến tính.
    \item \textbf{VAR (Vector AutoRegression):} Mở rộng ARIMA cho trường hợp đa biến, mô hình hóa tương quan giữa một số nút lân cận.
    \item \textbf{SVR và Gradient Boosting (XGBoost, LightGBM):} Ứng dụng các hàm nhân phi tuyến hoặc cây quyết định để hồi quy lưu lượng từ các đặc trưng trễ trong quá khứ.
\end{itemize}
Tuy nhiên, các phương pháp này gặp khó khăn khi mạng lưới mở rộng lên hàng trăm nút ($N = 608$) do số lượng tham số tăng nhanh, đồng thời thiếu cơ chế biểu diễn mối quan hệ topo không gian của mạng lưới đường bộ.

\subsection{Các mô hình học sâu tuần tự (RNN, LSTM, GRU)}
Với sự phát triển của học sâu, các kiến trúc mạng hồi quy tuần tự được áp dụng:
\begin{itemize}
    \item \textbf{LSTM và GRU:} Khắc phục hiện tượng triệt tiêu gradient của RNN truyền thống qua các cơ chế cổng điều tiết thông tin, ghi nhớ các phụ thuộc thời gian dài hạn hơn.
    \item \textbf{TCN (Temporal Convolutional Networks):} Sử dụng tích chập 1D giãn nở nhân quả, cho phép tính toán song song trên GPU và mở rộng trường tiếp nhận theo cấp số nhân.
\end{itemize}
Mặc dù vậy, các mô hình này chủ yếu xử lý độc lập từng nút hoặc làm phẳng toàn bộ mạng lưới thành một vector không có cấu trúc, bỏ qua thông tin hình học và liên kết vật lý giữa các tuyến đường.

\subsection{Mạng nơ-ron đồ thị không -- thời gian (STGNN)}
Để mô hình hóa đồng thời cả chiều không gian và thời gian, mạng nơ-ron đồ thị không -- thời gian (STGNN) đã ra đời và trở thành hướng nghiên cứu chủ đạo. Các kiến trúc STGNN tiêu biểu gồm:
\begin{enumerate}
    \item \textbf{Tích chập phổ trên đồ thị tĩnh:} Tiêu biểu là STGCN (Yu et al., 2018), sử dụng xấp xỉ ChebNet bậc $K=3$ trên ma trận Laplacian của đồ thị khoảng cách, kết hợp với các khối tích chập thời gian 1D Gated Linear Units (GLU). STGCN có tốc độ tính toán nhanh và khả năng định vị cục bộ tốt, nhưng phụ thuộc vào cấu trúc đồ thị tĩnh ban đầu.
    \item \textbf{Tích chập khuếch tán có hướng:} Tiêu biểu là DCRNN (Li et al., 2018), mô hình hóa dòng giao thông như một quá trình bước ngẫu nhiên có hướng trên đồ thị, kết hợp với tế bào GRU theo kiến trúc Encoder-Decoder. Hạn chế là tốc độ huấn luyện chậm do tính tuần tự của GRU.
    \item \textbf{Học ma trận kề thích ứng:} Tiêu biểu là GraphWaveNet (Wu et al., 2019), bổ sung ma trận kề tự học để nắm bắt tương quan tiềm ẩn giữa các nút xa nhau, kết hợp với tích chập thời gian giãn nở (TCN). Nhược điểm là ma trận kề tự học dạng dày làm tăng chi phí bộ nhớ khi số nút lớn và nhạy cảm với nhiễu.
    \item \textbf{Cơ chế chú ý không gian và thời gian phân tách:} Tiêu biểu là ASTGCN (Guo et al., 2019), tính toán ma trận chú ý động theo cả chiều không gian và thời gian. Tuy nhiên số lượng tham số lớn (1.53 triệu tham số) và chi phí tính toán tăng nhanh theo số lượng nút.
    \item \textbf{Đồ thị động và đồ thị meta:} Tiêu biểu là DSTAGNN (Lan et al., 2022) và MegaCRN (Jiang et al., 2023), xây dựng đồ thị tương quan động dựa trên khoảng cách Wasserstein hoặc phân tách biểu diễn đồ thị meta. Các mô hình này đạt độ chính xác tốt trên dữ liệu chuẩn nhưng dễ bị suy giảm khi đầu vào có nhiễu cảm biến lớn.
    \item \textbf{Mô hình Transformer không -- thời gian:} Các công trình như STAEformer (Liu et al., 2023) và iTransformer (Liu et al., 2024) ứng dụng cơ chế chú ý tự thân trên toàn chuỗi hoặc đảo ngược theo chiều biến số. Mặc dù có năng lực biểu diễn cao, các mô hình này đòi hỏi bộ nhớ VRAM đỉnh lớn (trên 11 GB với đồ thị 608 nút) và dễ bị suy thoái sai số ở các bước dự báo xa khi dữ liệu đầu vào mang nhiễu thực tế.
\end{enumerate}

\section{Hiện tượng lan truyền sai số trong hệ thống học máy nối tầng}
Trong các hệ thống học máy nối tầng, kết quả đầu ra của khâu nhận dạng đóng vai trò là dữ liệu đầu vào trực tiếp cho khâu dự báo:
\begin{equation}
    \mathbf{I} \xrightarrow{\quad f_{\theta_1} \quad} \hat{\mathbf{Y}} = \mathbf{Y}^* + \boldsymbol{\eta} \xrightarrow{\quad g_{\theta_2} \quad} \hat{\mathbf{Z}}
\end{equation}
Hiện tượng \textbf{lan truyền sai số (error propagation)} xuất hiện khi sai số nhận dạng $\boldsymbol{\eta}$ từ khâu $f_{\theta_1}$ không bị triệt tiêu mà lan truyền qua các tầng nơ-ron của khâu $g_{\theta_2}$, làm lệch lạc kết quả dự báo lưu lượng.

\subsection{Cơ chế tích lũy và khuếch đại sai số qua các bước thời gian}
Trong bài toán giao thông, hiện tượng này thể hiện ở hai khía cạnh:
\begin{enumerate}
    \item \textbf{Lan truyền sai số theo chiều không gian:} Phép tích chập đồ thị $g_\theta \star_{\mathcal{G}} \mathbf{x} = \sum_{k=0}^{K-1} \theta_k T_k(\tilde{\mathbf{L}}) \mathbf{x}$ khuếch tán thông tin trạng thái của một nút sang các nút lân cận $(K-1)$-hop. Nếu một camera bị lỗi nhận dạng cục bộ (ví dụ do mưa lớn hoặc lóa đèn pha), sai số này sẽ lan sang các nút xung quanh nếu mô hình không có cơ chế lọc phù hợp.
    \item \textbf{Tích lũy sai số theo chiều thời gian:} Khi dự báo đa bước từ $t+1$ đến $t+6$, sai lệch nhận dạng trong cửa sổ lịch sử 24 bước có thể làm biến dạng xu hướng ngắn hạn, dẫn đến việc ước lượng sai lệch tăng dần ở các bước nhìn xa.
\end{enumerate}

\subsection{Yêu cầu về độ bền vững với nhiễu nhận dạng}
Do đó, việc chỉ đánh giá mô hình trên dữ liệu sạch lý tưởng là chưa đủ đối với các hệ thống camera thực tế. Mô hình dự báo cần có \textbf{độ bền vững với nhiễu (noise robustness)}, tức năng lực tự động làm mượt và giảm thiểu ảnh hưởng của sai số đầu vào ngay trong cấu trúc nơ-ron. Đây là mục tiêu thiết kế then chốt của mô hình TA-STGCN.

\section{Đối sánh các kiến trúc dự báo giao thông tiêu biểu}
Để làm rõ bối cảnh nghiên cứu, Bảng \ref{tab:so_sanh_sota_toan_dien} tổng hợp các đặc tính kỹ thuật chính giữa các kiến trúc tiêu biểu hiện nay và mô hình đề xuất trong đề tài.

\begin{table}[htbp]
\centering
\footnotesize
\caption{Đối sánh đặc điểm các kiến trúc dự báo giao thông không -- thời gian tiêu biểu.}
\label{tab:so_sanh_sota_toan_dien}
\renewcommand{\arraystretch}{1.2}
\setlength{\tabcolsep}{3.5pt}
\begin{tabularx}{\textwidth}{l >{\raggedright\arraybackslash}p{2.6cm} >{\raggedright\arraybackslash}p{2.3cm} c >{\raggedright\arraybackslash}X c}
\toprule
\textbf{Mô hình} & \textbf{Cấu trúc không gian} & \textbf{Cơ chế thời gian} & \textbf{Cảm biến} & \textbf{Bối cảnh thực nghiệm} & \textbf{Phân loại} \\
\midrule
DCRNN (2018) & Khuếch tán đồ thị & GRU tuần tự & Vòng từ & Cao tốc ô tô (PeMS) & Không \\
STGCN (2018) & ChebNet ($K=3$) & Khối 1D GLU & Vòng từ & Cao tốc ô tô (PeMS) & Không \\
Graph WaveNet (2019) & Đồ thị thích ứng & Dilated TCN & Vòng từ & Cao tốc ô tô (PeMS) & Không \\
ASTGCN (2019) & ChebNet + Chú ý & Chú ý thời gian & Vòng từ & Cao tốc ô tô (PeMS) & Không \\
DSTAGNN (2022) & Đồ thị động đa kênh & Transformer thời gian & Vòng từ & Cao tốc ô tô (PeMS) & Không \\
MegaCRN (2023) & Meta-Graph thích ứng & GRU Encoder--Decoder & Vòng từ & Cao tốc ô tô (PeMS) & Không \\
STAEformer (2023) & Nhúng không gian & ST-Transformer & Vòng từ & Cao tốc ô tô (PeMS) & Không \\
iTransformer (2024) & Đảo chiều biến số & Inverted MLP & Chuỗi & Chuỗi đa biến chung & Không \\
\midrule
\textbf{TA-STGCN (Đề xuất)} & \textbf{ChebNet ($K=3$) thực} & \textbf{MHTA + GLU} & \textbf{Camera (608 nút)} & \textbf{Đô thị ($>80\%$ xe máy)} & \textbf{Có (2 lớp)} \\
\bottomrule
\end{tabularx}
\end{table}

\section{Khoảng trống nghiên cứu}
Từ khảo cứu trên, đề tài xác định ba khoảng trống nghiên cứu chính:
\begin{enumerate}
    \item \textbf{Thiếu tập dữ liệu chuẩn mực về giao thông hỗn hợp quy mô lớn:} Các tập dữ liệu phổ biến (PeMSD4, PeMSD8, METR-LA) đều thu thập từ cảm biến vòng từ trên cao tốc tại Mỹ, đặc trưng bởi dòng xe ô tô đơn nhất đi đúng làn. Chưa có tập dữ liệu chuẩn mực nào phục vụ bài toán dự báo toàn diện mạng lưới đô thị hỗn hợp với xe máy chiếm ưu thế trên quy mô toàn thành phố (608 camera).
    \item \textbf{Mất thông tin phân lớp phương tiện:} Hầu hết các mô hình hiện nay chỉ dự báo một kênh lưu lượng gộp chung hoặc tốc độ trung bình, làm mất đi đặc trưng động học riêng biệt của ô tô và xe máy.
    \item \textbf{Thiếu phân tích thực nghiệm về lan truyền sai số nhận dạng:} Các nghiên cứu trước đây thường giả định dữ liệu đầu vào là chính xác tuyệt đối. Vấn đề sai số nhận dạng từ camera ảnh hưởng như thế nào đến độ chính xác dự báo đa bước và cơ chế kiến trúc nào giúp nâng cao độ bền vững với nhiễu vẫn chưa được khảo sát đầy đủ.
\end{enumerate}
